\chapter{Einleitung}
\label{chap:einleitung}

\section{Einführung}
\label{sec:einleitung-einfuehrung}

Argumente bilden eine wichtige Grundlage menschlicher und maschineller Entscheidungsprozesse. Liegen widersprüchliche Informationen vor, muss bestimmt werden, 
welche Argumente trotz der bestehenden Konflikte gemeinsam vertreten werden können. Abstrakte Argumentationssysteme (engl.\ \emph{abstract argumentation frameworks}, kurz AFs) 
stellen hierfür ein einfaches und zugleich allgemeines Modell bereit. Sie wurden von \textcite{dung:1995:af} eingeführt und bilden Argumente als Knoten sowie Angriffe 
zwischen ihnen als gerichtete Kanten ab. Dabei wird bewusst vom Inhalt der Argumente und von der konkreten Ursache eines Konflikts abstrahiert. 

Für ein Argumentationssystem können unterschiedliche zulässige Bewertungen der enthaltenen Argumente bestehen. 
Ein Argument kann beispielsweise in einer vertretbaren Position akzeptiert und in einer anderen verworfen werden. 
Diese Mehrdeutigkeit bildet ab, dass eine Konfliktsituation mehrere gleichermaßen begründbare Standpunkte zulassen kann. 


Einen Ausgangspunkt für die Untersuchung solcher Zusammenhänge bietet die bedingte Unabhängigkeit. 
Sie stammt ursprünglich aus der Wahrscheinlichkeitstheorie und beschreibt, ob zwei Gruppen von Variablen noch Informationen übereinander liefern, 
nachdem eine dritte Gruppe beobachtet wurde. Grundlegende Eigenschaften solcher Unabhängigkeitsrelationen wurden unter anderem von
\textcite{pearl:paz:1986:graphoids} untersucht. Eine bekannte Anwendung bilden Bayes-Netze, in denen Unabhängigkeiten zwischen Variablen mithilfe
einer graphischen Darstellung sichtbar gemacht und für eine strukturierte Verarbeitung probabilistischer Informationen genutzt werden können \cite{pearl:2010}.

Die Grundidee bedingter Unabhängigkeit ist jedoch nicht auf Wahrscheinlichkeiten beschränkt.
Auch logische Wissensrepräsentationen können mehrere mögliche Zustände oder Interpretationen zulassen. 
Die von \textcite{darwiche:1997:logical-ci} entwickelte logische bedingte Unabhängigkeit beschreibt, 
ob Informationen aus verschiedenen Bereichen miteinander kombiniert werden können, sobald eine gemeinsame Hintergrundinformation festgelegt ist.

Diese Idee lässt sich auf abstrakte Argumentationssysteme übertragen. An die Stelle möglicher Interpretationen treten die 
zulässigen Bewertungen der Argumente. Zwei Argumentmengen können als bedingt unabhängig angesehen werden, 
wenn ihre möglichen Bewertungen nach Festlegung einer dritten Argumentmenge miteinander vereinbar bleiben. 
Bedingte Unabhängigkeit erfasst damit, ob zwei argumentative Teilbereiche getrennt voneinander betrachtet werden können.
 
Eine entsprechende Definition bedingter Unabhängigkeit für abstrakte Argumentationssysteme wurde von \textcite{rienstra:2020:independence-af} eingeführt. 
Das Erkennen solcher Unabhängigkeiten kann helfen, argumentative Zusammenhänge verständlicher zu machen und umfangreiche Probleme in kleinere Teilprobleme zu zerlegen.

Die exakte Entscheidung bedingter Unabhängigkeit ist für zentrale Auswertungsweisen abstrakter Argumentationssysteme 
jedoch komplexitätstheoretisch schwierig \cite{bluemel:2025:independence-aba}. 

Für die automatisierte Auswertung formaler Argumentationsprobleme wird häufig der Reduktionsansatz verwendet. 
Dabei wird die zu lösende Aufgabe in einen anderen Formalismus übersetzt, für den bereits leistungsfähige allgemeine 
Solver verfügbar sind. Aussagenlogische Kodierungen können beispielsweise von SAT-Solvern verarbeitet werden.
Quantifizierte boolesche Formeln erweitern diesen Ansatz und erlauben es, auch Beziehungen zwischen mehreren möglichen Bewertungen innerhalb einer Formel auszudrücken. 
Sie eignen sich daher besonders für die Formulierung bedingter Unabhängigkeit. Dass QBFs als einheitliches Zielformat für anspruchsvolle Aufgaben 
der formalen Argumentation eingesetzt werden können, zeigen bereits \textcite{egly:woltran:2006:reasoning-qbf}.

\section{Roadmap}
\label{sec:einleitung-roadmap}

Kapitel~2 führt die formalen Grundlagen der Arbeit ein. Dazu gehören abstrakte Argumentationssysteme, die betrachteten Auswertungsweisen, 
bedingte Unabhängigkeit, quantifizierte boolesche Formeln sowie die benötigten komplexitätstheoretischen Begriffe.

Kapitel~3 entwickelt QBF-Kodierungen für bedingte Unabhängigkeit unter der \emph{stable}-, \emph{complete}- und \emph{preferred}-Semantik. 
Diese Kodierungen werden anschließend verwendet, um ausgewählte ICCMA-Instanzen mithilfe eines QBF-Solvers auf Unabhängigkeit zu untersuchen. 
Die empirische Auswertung betrachtet, wie das Auftreten von Unabhängigkeit beziehungsweise Abhängigkeit mit der Größe der Konditionierungsmenge \(Z\) zusammenhängt. 
Zusätzlich wird untersucht, welchen Einfluss die Größe von \(Z\) auf die Solverlaufzeit besitzt.

Kapitel~4 analysiert die Komplexität der bedingten Unabhängigkeit für ausgewählte Graphklassen. Dabei wird untersucht, 
ob strukturelle Einschränkungen der Argumentationssysteme eine einfachere Entscheidung ermöglichen oder ob die hohe Komplexität des allgemeinen Problems erhalten bleibt.

\section{Motivation und Zielsetzung}
\label{sec:einleitung-motivation}

Ziel dieser Arbeit ist es, bedingte Unabhängigkeit in abstrakten Argumentationssystemen sowohl praktisch als auch theoretisch zu untersuchen. 
Dazu wird die Unabhängigkeitsentscheidung für die \emph{stable}-, \emph{complete}- und \emph{preferred}-Semantik als QBF formuliert.
Die Kodierung ermöglicht eine exakte und automatisierte Prüfung, ohne die zulässigen Bewertungen eines Argumentationssystems vorab vollständig aufzählen zu müssen. 
Dadurch wird eine empirische Analyse auf größeren Benchmarkinstanzen ermöglicht.

Von besonderem Interesse ist die Konditionierungsmenge \(Z\); sie bestimmt, welche Argumentbewertungen bereits als festgelegt gelten. 
Eine Veränderung ihrer Größe kann beeinflussen, ob zwei Argumentmengen unabhängig sind und wie aufwendig die entsprechende Solveranfrage ist. 
Die Untersuchung auf ICCMA-Instanzen soll zeigen, ob sich hierbei systematische Zusammenhänge beobachten lassen.

Ergänzend wird analysiert, welchen Einfluss die Graphstruktur eines Argumentationssystems auf die Schwierigkeit der Unabhängigkeitsentscheidung besitzt. 
Strukturelle Einschränkungen können die Auswertung eines Argumentationssystems vereinfachen. Dies muss jedoch nicht unmittelbar für bedingte Unabhängigkeit gelten, 
da diese die Beziehungen zwischen mehreren zulässigen Bewertungen betrachtet. Die komplexitätstheoretische Analyse soll daher feststellen,
 in welchen Graphklassen tatsächlich eine Vereinfachung erreicht wird und in welchen die Schwierigkeit des allgemeinen Falls bestehen bleibt.